\documentclass[12pt,a4paper]{article}
\usepackage[utf8]{inputenc}
\usepackage{amsmath}
\usepackage{amsfonts}
\usepackage{amssymb}
\usepackage{graphicx}
\usepackage{lmodern}
\usepackage[left=2cm,right=2cm,top=2cm,bottom=2cm]{geometry}
\usepackage{fontspec}
\setmainfont{Adobe Aldine}
\usepackage{csquotes}
\begin{document}

\section*{June 2026}

\subsection*{Monday, June 8}

I've been working on a new project of late, one that starts with a fairly big assumption: \textit{all} music from the 17th--21st centuries, save for serialist pieces, contains traces of contrapuntal foundations from 16th century music theories. Put in a less controversial way: there is at least one way to map out a path from what we hear in a piece of music written in 2026 to some form of contrapuntal theory or practice that existed (or was being developed) in the 16th century. My goal is to prove this by examining two types of musical categories: 

\begin{enumerate}
	\item \textbf{development} sections from pieces written in sonata form
	\item \enquote{\textbf{transitional}} works of music that exist on the edges of musical eras (e.g., \enquote{Classical,} \enquote{Romantic})
\end{enumerate}
\end{document}